\section{流程图：Diagrams as Simple Graphs}
\subsection{效果预览}
这个案例的最终效果如下
\begin{figure}[H]
    \centering
    \tikz [>={Stealth[round]}, black!50, text=black, thick,
        every new ->/.style = {shorten >=1pt},
        graphs/every graph/.style = {edges=rounded corners},
        skip loop/.style = {to path={-- ++(0,#1) -| (\tikztotarget)}},
        hv path/.style = {to path={-| (\tikztotarget)}},
        vh path/.style = {to path={|- (\tikztotarget)}},
        nonterminal/.style = {
            rectangle, minimum size=6mm, very thick, draw=red!50!black!50, top color=white,bottom color=red!50!black!20, font=\itshape, text height=1.5ex,text depth=.25ex},
        terminal/.style = {
            rounded rectangle, minimum size=6mm, very thick, draw=black!50, top color=white,bottom color=black!20, font=\ttfamily, text height=1.5ex, text depth=.25ex},
        shape = coordinate
    ]
    \graph [grow right sep, branch down=7mm, simple] {
        / -> unsigned integer[nonterminal] -- p1 -> "." [terminal] -- p2 -> digit[terminal] -- p3 -- p4 -- p5 -> E[terminal] -- q1 ->[vh path]
        {[nodes={yshift=7mm}]
            "+"[terminal], q2, "-"[terminal]
        } -> [hv path]
        q3 -- /unsigned integer [nonterminal] -- p6 -> /;
        p1 ->[skip loop=5mm] p4;
        p3 ->[skip loop=-5mm] p2;
        p5 ->[skip loop=-11mm] p6;
        q1 -- q2 -- q3; % make these edges plain
    };
    \caption{流程图效果预览}
\end{figure}

\subsection{节点设置}
\subsubsection{节点样式}

可以使用关键词 /.style 设置一个同一的样式
\begin{figure}[H]
    \centering
    \begin{minipage}{0.35\linewidth}
        \centering
        \begin{tikzpicture}[
                nonterminal/.style={
                rectangle,
                minimum size=6mm,
                very thick,
                draw=red!50!black!50, 
                top color=white, 
                bottom color=red!50!black!20, 
                font=\itshape
            }]
            \node [nonterminal] {unsigned integer};
        \end{tikzpicture}
    \end{minipage}
    \begin{minipage}{0.55\linewidth}
        \begin{lstlisting}[style = latex-side]
    \begin{tikzpicture}[
            nonterminal/.style={
            rectangle,
            minimum size=6mm,
            very thick,
            draw=red!50!black!50, 
            top color=white, 
            bottom color=red!50!black!20, 
            font=\itshape
        }]
        \node [nonterminal] {unsigned integer};
    \end{tikzpicture}
        \end{lstlisting}
    \end{minipage}
    \caption{基本节点样式}
\end{figure}

使用 rounded corner 进行圆角设置：
\begin{figure}[H]
    \centering
    \begin{minipage}{0.35\linewidth}
        \centering
        \begin{tikzpicture}[
            node distance = 5mm, terminal/.style = {
                rectangle,minimum size=6mm,rounded corners = 3mm,
                very thick, draw = black!50,
                top color = white, bottom color = black!20,
                font = \ttfamily
            }]
            \node (dot) [terminal] {.};
            \node (digit) [terminal,right=of dot] {digit};
            \node (E) [terminal,right=of digit] {E};
        \end{tikzpicture}
    \end{minipage}
    \begin{minipage}{0.55\linewidth}
        \begin{lstlisting}[style = latex-side]
    \begin{tikzpicture}[
        node distance = 5mm, terminal/.style = {
            rectangle,minimum size=6mm,rounded corners = 3mm,
            very thick, draw = black!50,
            top color = white, bottom color = black!20,
            font = \ttfamily
        }]
        \node (dot) [terminal] {.};
        \node (digit) [terminal,right=of dot] {digit};
        \node (E) [terminal,right=of digit] {E};
    \end{tikzpicture}
        \end{lstlisting}
    \end{minipage}
    \caption{节点样式：圆角}
\end{figure}

使用 shapes.misc 包的圆角设置:rounded rectangle
\begin{figure}[H]
    \centering
    \begin{minipage}{0.35\linewidth}
        \centering
        \usetikzlibrary {positioning}
        \begin{tikzpicture}[node distance=5mm,
            terminal/.style={
                rounded rectangle,
                minimum size=6mm,
                very thick,draw=black!50,
                top color=white,bottom color=black!20,
                font=\ttfamily
            }]
            \node (dot) [terminal] {.};
            \node (digit) [terminal,right=of dot] {digit};
            \node (E) [terminal,right=of digit] {E};
        \end{tikzpicture}
    \end{minipage}
    \begin{minipage}{0.55\linewidth}
        \begin{lstlisting}[style = latex-side]
    \usetikzlibrary {positioning}
    \begin{tikzpicture}[node distance=5mm,
        terminal/.style={
            rounded rectangle,
            minimum size=6mm,
            very thick,draw=black!50,
            top color=white,bottom color=black!20,
            font=\ttfamily
        }]
        \node (dot) [terminal] {.};
        \node (digit) [terminal,right=of dot] {digit};
        \node (E) [terminal,right=of digit] {E};
    \end{tikzpicture}
        \end{lstlisting}
    \end{minipage}
    \caption{rounded 圆角节点}
\end{figure}

在上述的几个例子中，我们发现 . 与 其他字符并没有对齐，. 看起来更像 ·。这是由于每个节点内容拥不同的高度与字母深度。可以使用 text height, text depth 来进行调整

\begin{figure}[H]
    \centering
    \begin{minipage}{0.35\linewidth}
        \centering
        \begin{tikzpicture}[
            node distance=5mm, text height=1.5ex, text depth = .25ex,
            terminal/.style={
                rounded rectangle,
                minimum size=6mm,
                very thick,draw=black!50,
                top color=white,bottom color=black!20,
                font=\ttfamily
            }]
            \node (dot) [terminal] {.};
            \node (digit) [terminal,right=of dot] {digit};
            \node (E) [terminal,right=of digit] {E};
        \end{tikzpicture}
    \end{minipage}
    \begin{minipage}{0.55\linewidth}
        \begin{lstlisting}[style = latex-side]
    \begin{tikzpicture}[
        node distance=5mm, text height=1.5ex, text depth = .25ex,
        terminal/.style={
            rounded rectangle,
            minimum size=6mm,
            very thick,draw=black!50,
            top color=white,bottom color=black!20,
            font=\ttfamily
        }]
        \node (dot) [terminal] {.};
        \node (digit) [terminal,right=of dot] {digit};
        \node (E) [terminal,right=of digit] {E};
    \end{tikzpicture}
        \end{lstlisting}
    \end{minipage}
    \caption{文字对齐}
\end{figure}

\subsubsection{节点位置}
使用位置选项 position options 设置节点相对位置
\begin{figure}[H]
    \centering
    \begin{minipage}{0.7\linewidth}
        \centering
        \begin{tikzpicture}[node distance=5mm and 5mm]
            \node (ui1) [nonterminal] {unsigned integer};
            \node (dot) [terminal,right=of ui1] {.};
            \node (digit) [terminal,right=of dot] {digit};
            \node (E) [terminal,right=of digit] {E};
            \node (plus) [terminal,above right=of E] {+};
            \node (minus) [terminal,below right=of E] {-};
            \node (ui2) [nonterminal,below right=of plus] {unsigned integer};
        \end{tikzpicture}
    \end{minipage}
    \begin{minipage}{0.6\linewidth}
        \begin{lstlisting}[style = latex-side]
    \begin{tikzpicture}[node distance=5mm and 5mm]
        \node (ui1) [nonterminal] {unsigned integer};
        \node (dot) [terminal,right=of ui1] {.};
        \node (digit) [terminal,right=of dot] {digit};
        \node (E) [terminal,right=of digit] {E};
        \node (plus) [terminal,above right=of E] {+};
        \node (minus) [terminal,below right=of E] {-};
        \node (ui2) [nonterminal,below right=of plus] {unsigned integer};
    \end{tikzpicture}
        \end{lstlisting}
    \end{minipage}
    \caption{节点位置}
\end{figure}

精确控制偏移距离的方法之一：使用 xshift yshift
\begin{figure}[H]
    \centering
    \begin{minipage}{0.35\linewidth}
        \centering
        \begin{tikzpicture}[scale = 1,node distance = 5mm and 5mm]
            \node (E) [terminal] {E};
            \node (plus) [terminal,above right=of E,xshift=5mm] {+};
            \node (minus) [terminal,below right=of E,xshift=5mm] {-};
            \node (ui2) [nonterminal,below right=of plus,xshift=5mm] {unsigned integer};
        \end{tikzpicture}
    \end{minipage}
    \begin{minipage}{0.55\linewidth}
        \begin{lstlisting}[style = latex-side]
    \begin{tikzpicture}[scale = 1,node distance = 5mm and 5mm]
        \node (E) [terminal] {E};
        \node (plus) [terminal,above right=of E,xshift=5mm] {+};
        \node (minus) [terminal,below right=of E,xshift=5mm] {-};
        \node (ui2) [nonterminal,below right=of plus,xshift=5mm] {unsigned integer};
    \end{tikzpicture}
        \end{lstlisting}
    \end{minipage}
    \caption{节点位置精确控制}
\end{figure}

\subsubsection{节点连线}
tikz 中的节点连接，尤其是非直线连接是一件非常复杂的事情。
\begin{figure}[H]
    \centering
    \begin{minipage}{0.35\linewidth}
        \centering
        \begin{tikzpicture}[node distance=5mm and 5mm]
            \node (dot) [terminal] {.};
            \node (digit) [terminal,right=of dot] {digit};
            \node (E) [terminal,right=of digit] {E};
            \path (dot) edge[->] (digit) % simple edges
                (digit) edge[->] (E);
            \draw [->]
                % start right of digit.east, that is, at the point that is the
                % linear combination of digit.east and the vector (2mm,0pt). We
                % use the ($ ... $) notation for computing linear combinations
                ($ (digit.east) + (2mm,0) $)
                % Now go down
                -- ++(0,-.5)
                % And back to the left of digit.west
                -| ($ (digit.west) - (2mm,0) $);
        \end{tikzpicture}
    \end{minipage}
    \begin{minipage}{0.55\linewidth}
        \begin{lstlisting}[style = latex-side]
    \begin{tikzpicture}[node distance=5mm and 5mm]
        \node (dot) [terminal] {.};
        \node (digit) [terminal,right=of dot] {digit};
        \node (E) [terminal,right=of digit] {E};
        \path (dot) edge[->] (digit) % simple edges
            (digit) edge[->] (E);
        \draw [->]
            % start right of digit.east, that is, at the point that is the linear combination of digit.east and the vector (2mm,0pt). We use the ($ ... $) notation for computing linear combinations
            ($ (digit.east) + (2mm,0) $)
            % Now go down
            -- ++(0,-.5)
            % And back to the left of digit.west
            -| ($ (digit.west) - (2mm,0) $);
    \end{tikzpicture}
        \end{lstlisting}
    \end{minipage}
    \caption{节点连线}
\end{figure}

在此基础上，tikz 引入了 edge 参数来简化折线。

\begin{figure}[H]
    \centering
    \begin{minipage}{0.35\linewidth}
        \centering
        \begin{tikzpicture}[
            node distance=5mm and 5mm,skip loop/.style={to path={-- ++(0,-.5) -| (\tikztotarget)}}]
            \node (dot) [terminal] {.};
            \node (digit) [terminal,right=of dot] {digit};
            \node (E) [terminal,right=of digit] {E};
            \path (dot) edge[->] (digit) % simple edges
                (digit) edge[->] (E)
                ($ (digit.east) + (2mm,0) $)
                        edge[->,skip loop] ($ (digit.west) - (2mm,0) $);
        \end{tikzpicture}
    \end{minipage}
    \begin{minipage}{0.55\linewidth}
        \begin{lstlisting}[style = latex-side]
    \begin{tikzpicture}[
        node distance=5mm and 5mm,skip loop/.style={to path={-- ++(0,-.5) -| (\tikztotarget)}}]
        \node (dot) [terminal] {.};
        \node (digit) [terminal,right=of dot] {digit};
        \node (E) [terminal,right=of digit] {E};
        \path (dot) edge[->] (digit) % simple edges
            (digit) edge[->] (E)
            ($ (digit.east) + (2mm,0) $)
                    edge[->,skip loop] ($ (digit.west) - (2mm,0) $);
    \end{tikzpicture}
        \end{lstlisting}
    \end{minipage}
    \caption{简化的节点连线}
\end{figure}

\subsection{矩阵布局}
矩阵式布局与 latex 自身的数学公式布局十分相似

\begin{figure}[H]
    \centering
    \begin{minipage}{0.7\linewidth}
        \centering
        \begin{tikzpicture}
            \matrix[row sep=1mm,column sep=5mm] {
                % First row:
                & & & & \node [terminal] {+}; & \\
                % Second row:
                \node [nonterminal] {unsigned integer}; &
                \node [terminal] {.}; &
                \node [terminal] {digit}; &
                \node [terminal] {E}; &&
                \node [nonterminal] {unsigned integer}; \\
                % Third row:
                & & & & \node [terminal] {-}; & \\
            };
        \end{tikzpicture}
    \end{minipage}
    \begin{minipage}{0.6\linewidth}
        \begin{lstlisting}[style = latex-side]
    \begin{tikzpicture}
        \matrix[row sep=1mm,column sep=5mm] {
            % First row:
            & & & & \node [terminal] {+}; & \\
            % Second row:
            \node [nonterminal] {unsigned integer}; &
            \node [terminal] {.}; &
            \node [terminal] {digit}; &
            \node [terminal] {E}; &&
            \node [nonterminal] {unsigned integer}; \\
            % Third row:
            & & & & \node [terminal] {-}; & \\
        };
    \end{tikzpicture}
        \end{lstlisting}
    \end{minipage}
    \caption{矩阵布局}
\end{figure}

使用了矩阵布局后，连线将变得非常简单

\begin{figure}[H]
    \centering
    \begin{minipage}{0.7\linewidth}
        \centering
        \begin{tikzpicture}
            [point/.style={circle,inner sep=0pt,minimum size=2pt,fill=red},skip loop/.style={to path={-- ++(0,#1) -| (\tikztotarget)}}]
            \matrix[row sep=1mm,column sep=2mm] {
                % First row:
                & & & & & & & & & & & \node (plus) [terminal] {+};\\
                % Second row:
                \node (p1) [point] {}; & \node (ui1) [nonterminal] {unsigned integer}; &
                \node (p2) [point] {}; & \node (dot) [terminal] {.}; &
                \node (p3) [point] {}; & \node (digit) [terminal] {digit}; &
                \node (p4) [point] {}; & \node (p5) [point] {}; &
                \node (p6) [point] {}; & \node (e) [terminal] {E}; &
                \node (p7) [point] {}; & &
                \node (p8) [point] {}; & \node (ui2) [nonterminal] {unsigned integer}; &
                \node (p9) [point] {}; & \node (p10) [point] {};\\
                % Third row:
                & & & & & & & & & & & \node (minus)[terminal] {-};\\
            };
            \path (p4) edge [->,skip loop=-5mm] (p3)
                    (p2) edge [->,skip loop=5mm] (p6);
        \end{tikzpicture}
    \end{minipage}
    \begin{minipage}{0.8\linewidth}
        \begin{lstlisting}[style = latex-side]
    \begin{tikzpicture}
        [point/.style={circle,inner sep=0pt,minimum size=2pt,fill=red},skip loop/.style={to path={-- ++(0,#1) -| (\tikztotarget)}}]
        \matrix[row sep=1mm,column sep=2mm] {
            % First row:
            & & & & & & & & & & & \node (plus) [terminal] {+};\\
            % Second row:
            \node (p1) [point] {}; & \node (ui1) [nonterminal] {unsigned integer}; &
            \node (p2) [point] {}; & \node (dot) [terminal] {.}; &
            \node (p3) [point] {}; & \node (digit) [terminal] {digit}; &
            \node (p4) [point] {}; & \node (p5) [point] {}; &
            \node (p6) [point] {}; & \node (e) [terminal] {E}; &
            \node (p7) [point] {}; & &
            \node (p8) [point] {}; & \node (ui2) [nonterminal] {unsigned integer}; &
            \node (p9) [point] {}; & \node (p10) [point] {};\\
            % Third row:
            & & & & & & & & & & & \node (minus)[terminal] {-};\\
        };
        \path (p4) edge [->,skip loop=-5mm] (p3)
                (p2) edge [->,skip loop=5mm] (p6);
    \end{tikzpicture}
        \end{lstlisting}
    \end{minipage}
    \caption{矩阵布局与连线}
\end{figure}

\subsection{连线}
\subsubsection{连接节点}
使用 graph 命令(path graph 的简写) 绘制连线 
\begin{figure}[H]
    \centering
    \begin{minipage}{0.7\linewidth}
        \centering
        \begin{tikzpicture}
            [point/.style={circle,inner sep=0pt,minimum size=2pt,fill=red},
                skip loop/.style={to path={-- ++(0,#1) -| (\tikztotarget)}},
                hv path/.style={to path={-| (\tikztotarget)}},
                vh path/.style={to path={|- (\tikztotarget)}}]
            \matrix[row sep=1mm,column sep=2mm] {
                % First row:
                & & & & & & & & & & & \node (plus) [terminal] {+};\\
                % Second row:
                \node (p1) [point] {}; & \node (ui1) [nonterminal] {unsigned integer}; &
                \node (p2) [point] {}; & \node (dot) [terminal] {.}; &
                \node (p3) [point] {}; & \node (digit) [terminal] {digit}; &
                \node (p4) [point] {}; & \node (p5) [point] {}; &
                \node (p6) [point] {}; & \node (e) [terminal] {E}; &
                \node (p7) [point] {}; & &
                \node (p8) [point] {}; & \node (ui2) [nonterminal] {unsigned integer}; &
                \node (p9) [point] {}; & \node (p10) [point] {};\\
                % Third row:
                & & & & & & & & & & & \node (minus)[terminal] {-};\\
            };
            \graph {
                (p1) -> (ui1) -- (p2) -> (dot) -- (p3) -> (digit) -- (p4)
                    -- (p5) -- (p6) -> (e) -- (p7) -- (p8) -> (ui2) -- (p9) -> (p10);
                (p4) -> [skip loop=-5mm] (p3);
                (p2) -> [skip loop=5mm] (p5);
                (p6) -> [skip loop=-11mm] (p9);
                (p7) -> [vh path] (plus) -> [hv path] (p8);
                (p7) -> [vh path] (minus) -> [hv path] (p8);
            };
        \end{tikzpicture}
    \end{minipage}
    \begin{minipage}{0.8\linewidth}
        \begin{lstlisting}[style = latex-side]
    \begin{tikzpicture}
        [point/.style={circle,inner sep=0pt,minimum size=2pt,fill=red},
            skip loop/.style={to path={-- ++(0,#1) -| (\tikztotarget)}},
            hv path/.style={to path={-| (\tikztotarget)}},
            vh path/.style={to path={|- (\tikztotarget)}}]
        \matrix[row sep=1mm,column sep=2mm] {
            ......
        };
        \graph {
            (p1) -> (ui1) -- (p2) -> (dot) -- (p3) -> (digit) -- (p4)
                -- (p5) -- (p6) -> (e) -- (p7) -- (p8) -> (ui2) -- (p9) -> (p10);
            (p4) -> [skip loop=-5mm] (p3);
            (p2) -> [skip loop=5mm] (p5);
            (p6) -> [skip loop=-11mm] (p9);
            (p7) -> [vh path] (plus) -> [hv path] (p8);
            (p7) -> [vh path] (minus) -> [hv path] (p8);
        };
    \end{tikzpicture}
        \end{lstlisting}
    \end{minipage}
    \caption{graph 连线}
\end{figure}

\subsubsection{连线样式}
\begin{figure}[H]
    \centering
    \begin{minipage}{0.7\linewidth}
        \centering
        \begin{tikzpicture}
            [>={Stealth[round]},thick,black!50,text=black,
                every new ->/.style={shorten >=1pt},
                graphs/every graph/.style={edges=rounded corners},
                point/.style={circle,inner sep=0pt,minimum size=2pt,fill=red},
                skip loop/.style={to path={-- ++(0,#1) -| (\tikztotarget)}},
                hv path/.style={to path={-| (\tikztotarget)}},
                vh path/.style={to path={|- (\tikztotarget)}}]
            \matrix[row sep=1mm,column sep=2mm] {
                % First row:
                & & & & & & & & & & & \node (plus) [terminal] {+};\\
                % Second row:
                \node (p1) [point] {}; & \node (ui1) [nonterminal] {unsigned integer}; &
                \node (p2) [point] {}; & \node (dot) [terminal] {.}; &
                \node (p3) [point] {}; & \node (digit) [terminal] {digit}; &
                \node (p4) [point] {}; & \node (p5) [point] {}; &
                \node (p6) [point] {}; & \node (e) [terminal] {E}; &
                \node (p7) [point] {}; & &
                \node (p8) [point] {}; & \node (ui2) [nonterminal] {unsigned integer}; &
                \node (p9) [point] {}; & \node (p10) [point] {};\\
                % Third row:
                & & & & & & & & & & & \node (minus)[terminal] {-};\\
            };
            \graph [use existing nodes] {
                p1 -> ui1 -- p2 -> dot -- p3 -> digit -- p4 -- p5 -- p6 -> e -- p7 -- p8 -> ui2 -- p9 -> p10;
                p4 ->[skip loop=-5mm] p3;
                p2 ->[skip loop=5mm] p5;
                p6 ->[skip loop=-11mm] p9;
                p7 ->[vh path] { plus, minus } -> [hv path] p8;
            };
        \end{tikzpicture}
    \end{minipage}
    \begin{minipage}{0.8\linewidth}
        \begin{lstlisting}[style = latex-side]
    \begin{tikzpicture}
        [>={Stealth[round]},thick,black!50,text=black,
            every new ->/.style={shorten >=1pt},
            graphs/every graph/.style={edges=rounded corners},
            point/.style={circle,inner sep=0pt,minimum size=2pt,fill=red},
            skip loop/.style={to path={-- ++(0,#1) -| (\tikztotarget)}},
            hv path/.style={to path={-| (\tikztotarget)}},
            vh path/.style={to path={|- (\tikztotarget)}}]
        \matrix[row sep=1mm,column sep=2mm] {
            ......
        };
        \graph [use existing nodes] {
            p1 -> ui1 -- p2 -> dot -- p3 -> digit -- p4 -- p5 -- p6 -> e -- p7 -- p8 -> ui2 -- p9 -> p10;
            p4 ->[skip loop=-5mm] p3;
            p2 ->[skip loop=5mm] p5;
            p6 ->[skip loop=-11mm] p9;
            p7 ->[vh path] { plus, minus } -> [hv path] p8;
        };
    \end{tikzpicture}
        \end{lstlisting}
    \end{minipage}
    \caption{连线样式}
\end{figure}

\subsection{Graph 指令}
使用 graph 节点可以快速绘制简单的流程图
\begin{figure}[H]
    \centering
    \begin{minipage}{0.6\linewidth}
        \centering
        \begin{tikzpicture}[scale = 1]
            \tikz \graph [grow right=2cm] {unsigned integer -> d -> digit -> E};
        \end{tikzpicture}
    \end{minipage}
    \begin{minipage}{0.6\linewidth}
        \begin{lstlisting}[style = latex-side]
    \tikz \graph [grow right=2cm] {unsigned integer -> d -> digit -> E};
        \end{lstlisting}
    \end{minipage}
    \caption{graph 指令绘图}
\end{figure}

声明 grow right sep 指明让节点在前一节点右方生成
\begin{figure}[H]
    \centering
    \begin{minipage}{0.7\linewidth}
        \centering
        \tikz \graph [grow right sep] {
            unsigned integer[nonterminal] -> "."[terminal] -> digit[terminal] -> E[terminal]
        };
    \end{minipage}
    \begin{minipage}{0.8\linewidth}
        \begin{lstlisting}[style = latex-side]
    \tikz \graph [grow right sep] {
        unsigned integer[nonterminal] -> "."[terminal] -> digit[terminal] -> E[terminal]
    };
        \end{lstlisting}
    \end{minipage}
    \caption{}
\end{figure}

创建分支
\begin{figure}[H]
    \centering
    \begin{minipage}{0.7\linewidth}
        \centering
        \tikz \graph [grow right sep] {
            unsigned integer [nonterminal] ->
            "." [terminal] ->
            digit [terminal] ->
            E [terminal] ->
            {
                "+" [terminal],
                "" [coordinate],
                "-" [terminal]
            } ->
            ui2/unsigned integer [nonterminal]
        };
    \end{minipage}
    \begin{minipage}{0.6\linewidth}
        \begin{lstlisting}[style = latex-side]
    \tikz \graph [grow right sep] {
        unsigned integer [nonterminal] ->
        "." [terminal] ->
        digit [terminal] ->
        E [terminal] ->
        {
            "+" [terminal],
            "" [coordinate],
            "-" [terminal]
        } ->
        ui2/unsigned integer [nonterminal]
    };
        \end{lstlisting}
    \end{minipage}
    \caption{}
\end{figure}
